\section{Further detail}
\label{app:detail}

The material below was moved out of the body for length. Nothing in it is altered.

\subsection{Checks, and what this module does not reach}

\emph{Supporting Section~\ref{sec:fiscal}, on the published figures the assembled rate is checked against and the limitation that runs against our own finding.}

Our shareholder layer at central values, \result{ShareholderLayerPct} percent, corroborates the
Congressional Research Service's text figure of about \result{CRSTextFigure} percent for the
effective capital gains rate on corporate profits. It does not reproduce that report's published
band of \result{CRSBandLow} to \result{CRSBandHigh} percent as printed, and should not: the band
already carries a taxable-share adjustment over domestically held stock. Rescaled to a common base
it is \result{CRSRescaledLow} to \result{CRSRescaledHigh} percent, and our figure sits inside it.
We report the correction rather than the original comparison, because as first written our own
check failed against our own construction.

The base construction treats foreign holders as bearing no tax at the owner level, though
dividends to them bear US withholding tax at treaty rates. That effect is sized in
Section~\ref{sec:fiscal}: applying a fifteen percent treaty rate to the foreign-held share raises
the assembled rate to \result{TauKWithholding} percent, and to \result{TauKJointWeakening} percent
together with the filers' own debt share, both still below the \result{RequiredAllGovEasy} percent
the condition requires. The treaty rate is a stated assumption, not a measurement, and what stays
open is the effective rate actually borne, which varies by treaty and with the split between
dividends and capital gains realized by foreign holders. The debt share is an economy-wide stock ratio, \result{DebtShareLow} to
\result{DebtShareHigh} percent, standing in for a marginal flow; the evidence in
Section~\ref{sec:bust} points below that range, and because the debt-financed term is negative, a
lower debt share would raise the assembled rate.

\subsection{The pass-through decomposition}

\emph{Supporting Section~\ref{sec:fiscal}, on the taxable capital income generated per displaced wage dollar.}

Equation~\eqref{eq:passthrough} in Section~\ref{sec:fiscal} writes the condition with the
pass-through coefficient $g$ explicit; the headline sets $g = 1$ because $g$ cannot be sourced.
What follows is the sensitivity behind that choice, and its only purpose is to show the
direction: incomplete pass-through can only widen the gap on the marginal rate, never close it.

A displaced wage dollar splits into the cost of the AI capital and services that replace the work,
a share $c$, and the surplus the adopting firm keeps, $1 - c$. The second is the adopting firm's
capital income. Only the first is revenue to the suppliers of AI capital, and only there does our
decomposition of a dollar of AI capital spending apply: \result{CapexDomesticLabor} percent
domestic labor, which is already inside $\tau_\ell (1 - R)$ and must not be counted twice,
\result{CapexImported} percent imported and never entering the US tax base, and
\result{CapexDomesticSurplus} percent domestic producer surplus. So
$g = (1 - c) + c \, s$, with $s$ that surplus share. Taking $c = \result{CapitalCostShare}$
percent gives $g = \result{GCentral}$ and a required rate of \result{RequiredPTEasy} to
\result{RequiredPTHarder} percent, against \result{RequiredEasy} to \result{RequiredHard} at
$g = 1$; across the swept range of $c$, $g$ runs from \result{GLow} to \result{GHigh}.

A second channel is reported separately because it is a different mechanism. If a share of the
adopting firm's retained surplus is competed away to consumers in lower prices, it is never taxed
as capital income. At \result{PriceThrough} percent of the retained surplus, $g$ falls to
\result{GWithPrice} and the required rate rises to \result{RequiredPriceEasy} to
\result{RequiredPriceHarder} percent.

Three statements about this sensitivity, since it is the weakest thing in the section. It is
scenario arithmetic: $c$ has no verified estimate here or, so far as we located, anywhere, and the
shares behind $s$ are swept ranges rather than sourced values. It does not double count the
profit shifting already inside the assembled rate: the import share removed here is payment that
never enters the US base, while shifted profit does enter it and is taxed at the lower rate the
assembly already applies, and the one case where the two could touch, an intra-firm import priced
as a transfer, would make $g$ too low and so runs against this correction. And it changes what can
be said about the economy-wide rate only conditionally: at the central case of the base
construction, an economy-wide rate of twenty to twenty-two percent still clears the required band;
it stops clearing once a price pass-through term is added. That is a reason to report the two
terms separately, not to fold either into the headline.

\subsection{The transition path under expensing}

\emph{Supporting Section~\ref{sec:fiscal}, on revenue during an investment boom.}

While investment grows, current deductions exceed the tax on income from a capital stock built
when investment was smaller, so net revenue from the sector is below the steady-state figure and
can be negative. An illustrative path, using the tax parameters of Section~\ref{sec:fiscal} and a
geometric investment series that is chosen rather than estimated, makes the size of that visible:
with investment growing at \result{TransitionGrowthRate} percent a year, net entity-level tax is
negative at \result{TransitionYearFive} per unit of investment in year five and still negative in
year twenty-five, while at zero growth it turns positive in year
\result{TransitionPositiveYear}. That path illustrates a mechanism and is not a forecast of
revenue. The mechanism itself is not in doubt: during an investment boom the revenue actually
collected is lower than the assembled rate implies, not higher.

\subsection{The financing mix inside the measured tier}

\emph{Supporting Section~\ref{sec:bust}, on how the observed AI investment is funded.}

The nine filers report \result{FilerCapexBn} billion dollars of capital spending against
\result{FilerOCFBn} billion of operating cash flow, a self-funding ratio of
\result{SelfFunding}, and a debt-financed share of capital spending of
\result{DebtFinancedShare} percent on those books, which is a lower bound because financing held
off those balance sheets does not appear in it. Of the bank commitments in the second tier,
roughly \result{AIDrawnBn} billion appears drawn \citep{CohenKillenLau2026}.

\subsection{Reading the lever ranking}

\emph{Supporting Section~\ref{sec:fiscal}, on the table of levers.}

One line of Table~\ref{tab:levers} is easy to read backwards. The taxable-holder share appears
last, with a swing of \result{SwingTheta}, because it is pinned to a narrow sourced range and not
because it does not matter: it is the largest single determinant of the level of the assembled
rate and the smallest remaining source of uncertainty in it.
